\documentclass[proposal,english,palatino]{AIGpaper}

%%%% Package Imports %%%%%%%%%%%%%%%%%%%%%%%%%%%%%%%%%%%%%%%%%%%%%%%%%%%%%%%%%%%%%%%%%
\usepackage{graphicx}					    % enhanced support for graphics
\usepackage{tabularx}				      	% more flexible tabular
\usepackage{amsfonts}					    % math fonts
\usepackage{amssymb}					    % math symbols
\usepackage{amsmath}					    % overall enhancements to math environment

%%%% optional packages
\usepackage{tikz}                           % creating graphs and other structures
\usetikzlibrary{arrows,positioning}
\tikzset{
    %Define standard arrow tip
    >=stealth',
    %Define style for argument
    args/.style={circle, minimum size=0.9cm,draw=black, thick,fill=white},
}

%%%% Author and Title Information %%%%%%%%%%%%%%%%%%%%%%%%%%%%%%%%%%%%%%%%%%%%%%%%%%%
\input{meta.tex}
\author{\authorname}

\title{Inconsistency Measurement in Classical Planning Problems}

%%%% Abstract %%%%%%%%%%%%%%%%%%%%%%%%%%%%%%%%%%%%%%%%%%%%%%%%%%%%%%%%%%%%%%%%%%%%%%
\englishabstract{
    Classical planning problems become inconsistent when no valid action sequence can transform the initial state into a goal state. While inconsistency measurement has been well-developed in knowledge representation, its application to planning domains has received limited attention. This thesis aims to explore how established inconsistency measures can be adapted to classical planning problems in STRIPS and PDDL, with a focus on developing Answer Set Programming (ASP)-based algorithms for their computation. The work will investigate the practical utility of these measures for problem diagnosis and debugging. The goal is to provide planners with better tools for understanding why problems fail, potentially supporting more systematic approaches to constraint relaxation and problem reformulation.
}

\begin{document}

\maketitle % prints title and author information, as well as the abstract 

% ===================== Beginning of the actual text section =====================

\section{Introduction and Motivation}

Classical planning problems in artificial intelligence involve finding action sequences that transform an initial state into a goal state. When no valid action sequence exists, the problem is considered unsolvable. Most planning approaches focus on solvable instances, and when faced with unsolvable problems, they typically just report failure without much additional insight.

Inconsistency measurement is an established field in knowledge representation that quantifies the degree of inconsistency in logical knowledge bases. It provides ways to understand conflicts within sets of propositions and has been applied to various domains. However, its application to planning appears to be limited.

The motivation for this work comes from observing that many unsolvable planning problems seem to contain conflicting constraints or goals. For example, a problem might require an object to be in two different locations simultaneously, or actions might have effects that make certain goal combinations impossible. When planning systems encounter such conflicts, they usually just return "unsolvable" without indicating what went wrong or how the problem might be fixed.

It seems reasonable that quantifying the degree of inconsistency could help identify conflict sources and guide problem reformulation. This could be particularly useful when working with complex planning domains or automatically generated problems, where manually debugging failures can be tedious.

However, adapting inconsistency measures from knowledge representation to planning is not straightforward. Planning problems have a temporal structure and involve sequences of actions and state changes, which differs significantly from static logical knowledge bases. The types of conflicts that arise in planning may also be quite different from those in propositional logic. This suggests that existing measures will need careful adaptation to capture the specific characteristics of planning problems.

\section{Technical Foundations and Related Work}

\subsection{Inconsistency Measurement in Knowledge Representation}

Inconsistency measurement is a research area in knowledge representation that deals with quantifying the degree of inconsistency in logical knowledge bases. Several types of measures have been developed:

\textbf{Distance-based measures} calculate the minimal changes needed to make a knowledge base consistent. For example, the Max-Distance and Sum-Distance measures compute the Hamming distance between inconsistent models and their nearest consistent versions \cite{Kuhlmann2025}.

\textbf{Hitting-set-based measures} work by identifying minimal inconsistent subsets and quantifying their overlaps. The Hitting-Set and Hit-Distance measures fall into this category and appear to be computationally reasonable \cite{Kuhlmann2025}.

\textbf{Contention-based measures} count conflicting formulas or statements. The Contention measure seems to be one of the most studied in the literature \cite{Kuhlmann2025}.

Recent work suggests that Answer Set Programming (ASP)-based algorithms for computing inconsistency measures can sometimes outperform SAT-based approaches, especially for more complex measures \cite{Kuhlmann2025}. This might be relevant for planning applications, since ASP's declarative nature could be well-suited for modeling planning scenarios.

Additionally, Ulbricht et al. \cite{Ulbricht2016} developed inconsistency measures specifically for Answer Set Programs, showing that classical monotonicity properties don't hold in non-monotonic logics. This could be relevant when adapting measures to planning, where non-monotonic aspects might also play a role.

\subsection{Classical Planning}

Classical planning works with a state-based formalism where states are described by sets of facts, and actions are defined through preconditions and effects. The STRIPS (Stanford Research Institute Problem Solver) representation is foundational for many modern planning systems \cite{Russell2016}.

A planning problem becomes unsolvable when no sequence of applicable actions can transform the initial state into a goal state. The computational complexity of determining plan existence is PSPACE-complete \cite{Russell2016}, which has motivated research into efficient heuristics and algorithms.

\textbf{PDDL (Planning Domain Definition Language)} is the standard for describing planning problems and is used in international planning competitions. Adapting inconsistency measures to work with PDDL representations would be important for practical applicability.

Modern planning systems use various search strategies and heuristics, but when they encounter unsolvable problems, they typically just report failure without much analysis of why the problem failed or what might be done about it.

\subsection{Planning and Inconsistency: Foundation Work}

Eriksson et al. \cite{Eriksson2018} established the theoretical connection between unsolvable planning problems and inconsistency. Their work focuses on proof systems for demonstrating unsolvability in planning, particularly examining why certain planning problems cannot be solved using STRIPS-based approaches.

However, their work is primarily concerned with qualitative aspects of unsolvability rather than quantitative measures of inconsistency degrees. While they provide important groundwork by formalizing the relationship between planning failures and logical inconsistency, the question of how to measure and compare different degrees of inconsistency in planning problems remains open.

This gap suggests there might be value in adapting established inconsistency measures from knowledge representation to the planning domain, building on the theoretical foundations they established.

\section{Problem Statement and Research Objectives}

This thesis aims to explore how inconsistency measures from knowledge representation can be adapted to classical planning and whether they provide practical value for understanding planning failures.

\subsection{Problem Statement}

Currently, when planning systems encounter unsolvable problems, they provide limited diagnostic information. This creates several challenges:

\begin{itemize}
    \item Users have little insight into why problems fail
    \item There's no systematic way to analyze the structure of problematic planning instances
    \item Support for reformulating problems is limited
    \item Different types of planning failures are not distinguished quantitatively
\end{itemize}

\subsection{Research Questions and Objectives}

This work will address the central research question:

\textbf{How can inconsistency measures from knowledge representation be effectively adapted to classical planning to provide practical diagnostic value for unsolvable planning problems?}

This main question encompasses several sub-objectives:

\textbf{Theoretical Adaptation:} Investigate how to systematically transfer representative inconsistency measures (contention-based, hitting-set-based, and distance-based) to classical planning, developing translation mechanisms between PDDL representations and inconsistency measure frameworks while examining which theoretical properties can be preserved.

\textbf{Algorithmic Development:} Explore ASP-based algorithms for computing the adap-ted measures, implement prototypes that work with standard planning representations, and investigate optimization techniques where computational performance becomes critical.

\textbf{Empirical Evaluation:} Evaluate the measures on planning benchmarks from international competitions, analyze computational costs and scalability, and assess whether the measures provide interpretable and practically useful information for planners.

\textbf{Practical Applications:} Explore diagnostic tools for planning problems, investigate applications to problem debugging and constraint relaxation, and consider integration possibilities with existing planning tools and workflows.

\section{Planned Structure of the Master's Thesis}

The master's thesis will be structured in eight chapters. After the introduction (Chapter 1), the foundations will cover classical planning formalisms, inconsistency measurement theory, and Answer Set Programming (Chapter 2). Related work will review unsolvable planning problems, state-of-the-art inconsistency measures, and computational approaches (Chapter 3). The theoretical core will develop formal translations of selected measures and analyze their properties and complexity (Chapter 4). The implementation phase will focus on ASP-based algorithms and PDDL integration (Chapter 5). Empirical evaluation will include experimental setup, benchmark evaluation, and performance analysis (Chapter 6). Applications and case studies will explore diagnostic tools and practical scenarios (Chapter 7). The conclusion will summarize results, discuss limitations, and outline future work (Chapter 8).

\section{Timeline}

The planned timeline for the master's thesis spans 6 months with overlapping phases. Months 1-2 will focus on literature review and foundations, including in-depth study of inconsistency measures, familiarization with classical planning theory and tools, and completion of the theoretical groundwork (Chapters 2-3). Months 2-3 will emphasize theoretical development through formal adaptation of selected measures, proofs of theoretical properties, and complexity analysis (Chapter 4). Months 3-4 will concentrate on implementation, developing ASP-based algorithms, PDDL integration, and initial prototype testing (Chapter 5). Months 4-5 will involve systematic evaluation including benchmark testing, performance analysis, and comparative studies (Chapter 6). Finally, months 5-6 will focus on applications and finalization through case studies, diagnostic tool development, and completing the writing and revision process (Chapters 1, 7, 8).

% References
\addreferences

\end{document}
